\chapter{The remaining case}\label{ch:remaining}

This chapter attempts the general case. Fix a Weil family
$\pi\colon\mathcal A\to\mathrm{Sh}$ of dimension $2n\ge4$, as in
Chapter~\ref{ch:weil}, and a nonzero Weil class $\omega$, extended to a flat
global section by Proposition~\ref{prop:flat}. By
Corollary~\ref{cor:oneclass}, the Hodge conjecture for a member with full
Hodge group is the algebraicity of $\omega$ there. The attempt follows the
most direct route: take a member where $\omega$ is algebraic and move the
cycle to every other member. We prove what that route gives and show
exactly where it stops. Then we look at the problem as a case of the
variational Hodge conjecture and prove three results that narrow where a
proof could come from.

Much of the material in the first sections first appeared in my preprint
\cite{BhH8}. Everything used is proved here.

\section{The algebraic locus}

\begin{definition}
The \emph{algebraic locus} is the set $\Sigma\subseteq\mathrm{Sh}$ of points
$t$ at which $\omega(t)$ is the class of an algebraic cycle with rational
coefficients on $A_{t}$.
\end{definition}

The goal is $\Sigma=\mathrm{Sh}$.

\begin{lemma}\label{lem:countablebook}
$\Sigma$ is a countable union of closed algebraic subsets of $\mathrm{Sh}$.
\end{lemma}

\begin{proof}
Fix a relatively ample line bundle on $\mathcal A/\mathrm{Sh}$. For each Hilbert
polynomial $P$ the relative Hilbert scheme $\Hilb_{P}(\mathcal A/\mathrm{Sh})$ is
projective over $\mathrm{Sh}$ and has finitely many irreducible components. Take
finitely many components $H_{1},\dots,H_{r}$, possibly from different
Hilbert schemes, and rational numbers $q_{1},\dots,q_{r}$. Over the fibre
product $H=H_{1}\times_{\mathrm{Sh}}\dots\times_{\mathrm{Sh}}H_{r}$ we have the
class $\sum_{i}q_{i}[Z_{i}]$ in the fibre $H^{2n}(A_{t},\QQ)$. Classes of
the fibres of a flat family of subschemes are locally constant. So on each
connected component of $H$ the class $\sum q_{i}[Z_{i}]$ either equals
$\omega(t)$ everywhere or nowhere. The union of the components where it
equals $\omega(t)$ maps properly to $\mathrm{Sh}$, so its image is closed and
algebraic. There are countably many choices of polynomials, components
and coefficients, and every point of $\Sigma$ lies in one of these images.
\end{proof}

\begin{lemma}[Dichotomy]\label{lem:dichotomybook}
Either $\Sigma=\mathrm{Sh}$, or $\Sigma$ is meagre: it lies in a countable
union of closed subsets with empty interior, and its complement is dense.
\end{lemma}

\begin{proof}
$\mathrm{Sh}$ is irreducible. If one of the closed sets of
Lemma~\ref{lem:countablebook} is all of $\mathrm{Sh}$, then
$\Sigma=\mathrm{Sh}$. Otherwise each of them is a proper closed algebraic
subset, which has empty interior, and Baire's theorem on the complex
manifold $\mathrm{Sh}$ gives the rest.
\end{proof}

By Section~\ref{sec:cmpoint}, every Weil family has a member $s_{0}$ at
which $\omega$ is a polynomial in divisor classes, and it can be chosen
with complex multiplication \cite{BhH8}. By Proposition~\ref{prop:densebook}
the isogeny orbit of $s_{0}$ lies in $\Sigma$ and is dense. So $\Sigma$ is
dense, and the problem is entirely about which side of the dichotomy we are
on. One stratum of Lemma~\ref{lem:countablebook} equal to all of
$\mathrm{Sh}$ would settle it.

There is a catch, and it is the reason the problem is hard. Recall that
$\mathrm{Sh}^{\circ}$ is the set of points where the Hodge group is all of
$\SU(V,H)$. Its complement is meagre (Proposition~\ref{prop:generalsu}).
Every point isogenous to a CM point is a CM point, with a torus for its
Hodge group, so the whole orbit of $s_{0}$ lies \emph{outside}
$\mathrm{Sh}^{\circ}$. To close the locus we need at least one point of
$\Sigma\cap\mathrm{Sh}^{\circ}$. For $K=\QQ(\sqrt{-3})$ and a split form,
Schoen's Prym varieties supply such points (Section~\ref{sec:schoen}). For
other fields and discriminants no such point is known when $n\ge3$.

\section{The first route: strata with big monodromy}

Suppose that we had found one point $t\in\Sigma\cap\mathrm{Sh}^{\circ}$, and let
$S$ be an irreducible stratum of Lemma~\ref{lem:countablebook} through it. The
first thing to ask is what $S$ can look like. The following lemma says that
$S$ cannot be a small, degenerate piece of the family.

\begin{theorem}\label{thm:monodromybook}
Let $S\subseteq\mathrm{Sh}$ be an irreducible closed subvariety of positive
dimension that contains a point of $\mathrm{Sh}^{\circ}$. Then the connected
algebraic monodromy group of the local system $R^{1}\pi_{*}\QQ$ restricted
to the smooth locus of $S$ is the whole group $\SU(V,H)$.
\end{theorem}

\begin{proof}
Let $S^{\mathrm{reg}}$ be the smooth locus of $S$. It is connected, since
$S$ is irreducible. The Hodge group of a very general point of $S$ contains
the Hodge group of every point of $S$, so it contains $\SU(V,H)$. On the
other hand every Hodge group in the family lies in $\SU(V,H)$, by
\eqref{eq:hginsu}.
So the generic Hodge group on $S$ is $\SU(V,H)$.

By Andr\'e's theorem \cite{And92}, the connected algebraic monodromy group
$M$ of a polarisable variation of Hodge structure over a smooth connected
base is a normal subgroup of the derived group of the generic
Mumford--Tate group. Here that derived group is $\SU(V,H)$, which is an
outer form of $\SL_{2n}$ over $\QQ$, and so it is almost simple. Its
connected normal subgroups are $1$ and itself.

If $M=1$, the monodromy is finite. After a finite \'etale cover of
$S^{\mathrm{reg}}$ the local system becomes trivial, and the period map
lifts to a holomorphic map from that cover to $D$. Now $D$ is a bounded
domain, so the coordinate functions of the lift are bounded holomorphic
functions. Pass to a resolution and a smooth compactification of the
cover. By the Riemann extension theorem the bounded functions extend to the
compactification, which is compact and connected, so they are constant.
Then the period map is constant on $S^{\mathrm{reg}}$. But the period map of
the universal family over $\mathrm{Sh}$ is the identity of $\Gamma\backslash D$, so
$S$ is a point, which contradicts $\dim S>0$. Hence $M=\SU(V,H)$.
\end{proof}

So a positive-dimensional stratum through a general point carries all the
monodromy the family has. It is natural to hope that a cycle family with
this much monodromy can be moved off $S$. That hope is exactly the
variational Hodge conjecture for this family, and here is why the
monodromy does not help on its own. The class $\omega$ is invariant under
all of $\SU(V,H)$, so the monodromy never moves $\omega$, and it gives no
constraint that would force the stratum to grow. Big monodromy rules out
degenerate strata. It does not produce new cycles.

\section{The second route: deforming one cycle}

The direct way to move a cycle is deformation theory. Let $Z\subset A_{t}$
be a local complete intersection of codimension $n$. Bloch \cite{Blo72}
defined the \emph{semiregularity map}
\[
  \sigma_{Z}\colon H^{1}(Z,N_{Z})\longrightarrow H^{n+1}(A_{t},\Omega^{n-1}),
\]
and proved the following. If $\sigma_{Z}$ is injective, then $Z$ deforms
along every deformation of $A_{t}$ along which its class stays of Hodge type.
Buchweitz and Flenner \cite{BF03} extended the definition and the theorem
to perfect complexes, and Pridham \cite{Pri24} gave a general proof.

\begin{theorem}\label{thm:semiregbook}
Suppose that at some point $t\in\mathrm{Sh}$ there is a local complete intersection
$Z\subset A_{t}$ of codimension $n$ whose class is
$\cl(Z)=\alpha+c\,\eta^{n}$ with $\alpha$ a nonzero Weil class and
$c\in\QQ$, and whose semiregularity map is injective. Then $\Sigma=\mathrm{Sh}$.
\end{theorem}

\begin{proof}
Along the universal family over a small neighbourhood $U$ of $t$, the flat
transport of $\cl(Z)$ is $\alpha+c\eta^{n}$, and both terms are Hodge
classes at every point, because the Weil plane and $\eta$ are. So the class
of $Z$ stays of Hodge type along every direction of $U$. By Bloch's theorem,
$Z$ deforms to every order along every curve germ in $U$. Then the image of
its Hilbert scheme component contains the formal neighbourhood of $t$. Since
that image is a closed analytic subset of $U$, it contains a neighbourhood
of $t$. So one stratum of Lemma~\ref{lem:countablebook} has nonempty interior,
and it is all of $\mathrm{Sh}$ by irreducibility. At every point $s$ we then have an
algebraic cycle of class $\alpha(s)+c\eta^{n}$, so $\alpha(s)$ is
algebraic.

It remains to pass from $\alpha$ to $\omega$. The field $K$ acts on the
Weil plane: an element $\tau\in K^{\times}$ acts on $\bigwedge^{2n}_{K}V$ by
multiplication by $\tau^{2n}$. A positive integer multiple $m\tau$ is a
genuine endomorphism of $A_{s}$. Its pullback, divided by $m^{2n}$, is
$\tau^{*}$ on $H^{2n}$, so $\tau^{*}$ carries algebraic classes to
algebraic classes. Choose $\tau$ whose argument is not a rational multiple of $\pi$, for
example $\tau=a+\sqrt{-d}$ with $a$ a positive integer chosen so that
$\tau/\bar\tau$ is not a root of unity. Then $\tau^{2n}\notin\QQ$, and the classes
$\alpha$ and $\tau^{*}\alpha$ span the Weil plane over $\QQ$, so $\omega(s)$
is algebraic.
\end{proof}

The same holds with $Z$ replaced by a perfect complex $E$ on $A_{t}$, with
the Chern character in place of the cycle class, by the theorem of
Buchweitz--Flenner and Pridham. This is the route Markman followed for
sixfolds of split Weil type, using secant sheaves and an equivariant form
of semiregularity \cite{Mar25b}.

\section{Where the attempt stops}\label{sec:stopsbook}

Both routes need an input that we do not have.

\subsection*{The first route needs a point of $\Sigma$ with full Hodge group.}
Outside the families where the conjecture is already known and Schoen's
Prym locus for $\QQ(\sqrt{-3})$, every point at which $\omega$ is known to
be algebraic is a CM point or lies on a proper special subvariety. These
points are excluded from $\mathrm{Sh}^{\circ}$ by definition.
Theorem~\ref{thm:monodromybook} tells us what a stratum through a point of
$\mathrm{Sh}^{\circ}$ looks like, but it gives no way to find one. Where one is
known, Schoen's Prym locus, it has dimension $3n$, while $\mathrm{Sh}$ has
dimension $n^{2}$. Nothing in the first route makes such a stratum grow.

\subsection*{The second route needs a semiregular object, and the obvious
ones are not.}
The first thing to try is an object whose Chern character is exactly a
Weil class, with no other component. Voisin's tori rule this out.

\begin{theorem}\label{thm:purefalse}
Let $n\ge2$, let $t\in\mathrm{Sh}$, and let $F$ be a perfect complex on
$A_{t}$ with $\operatorname{ch}(F)=N\omega(t)$ for some $N\ne0$, possibly
plus components in degrees $0$ and $4n$. Then the semiregularity map of $F$
is not injective.
\end{theorem}

\begin{proof}[Sketch of proof]
Suppose it is injective. The class $N\omega$ is of Hodge type not only on
the members of the Weil family but on every Weil torus of
Section~\ref{sec:voisintori}, which form a connected complex manifold
$\mathcal S$ of dimension $2n^{2}$ containing $\mathrm{Sh}$ locally. So the
theorem of Buchweitz--Flenner and Pridham \cite{BF03,Pri24} kills every
obstruction to deforming $F$ along an Artinian deformation of $A_{t}$
inside $\mathcal S$. Kuranishi's method then gives an analytic family of
complexes of vector bundles over a neighbourhood of $t$ in $\mathcal S$,
deforming $F$. At a very general Weil torus $T$ in that neighbourhood,
every coherent sheaf has $\operatorname{ch}_{k}=0$ for $0<k<2n$, by the
argument after Theorem~\ref{thm:voisin}. So $\operatorname{ch}_{n}$ of the
deformed complex vanishes on $T$, while it is the flat transport of
$N\omega\ne0$. This is a contradiction. The details are written out in
\cite{BhH8}.
\end{proof}

So a semiregular object has to carry a polarisation term, as in
Theorem~\ref{thm:semiregbook}, and that term must not stay of Hodge type on
the Weil tori, or the same argument applies. With a polarisation term the
criterion survives, but it asks for an object with a precise and large
amount of self-extension. For $n\ge3$ and a general imaginary quadratic
field no such object is known, and the natural constructions that provably
fail are listed in \cite{BhH8}.

\subsection*{Neither gap is a technicality.}
If $\Sigma=\mathrm{Sh}$, then $\Sigma$ meets $\mathrm{Sh}^{\circ}$, so the first gap is
a consequence of the goal as well as a step towards it. And closing the
locus is equivalent to one case of the Lefschetz standard conjecture: the
case of the total space of the family over a projective curve through
$s_{0}$ and a very general point \cite{BhH8}. These are the variational Hodge
conjecture and the standard conjectures in disguise. They are not small residuals.

\section{The variational form of the problem}\label{sec:vhcbook}

The question whether $\Sigma=\mathrm{Sh}$ is one instance of the variational Hodge
conjecture, and it helps to look at it that way. Let $f\colon X\to S$ be a
smooth projective morphism onto a smooth connected variety, and let
$\alpha$ be a global section of $R^{2p}f_{*}\QQ$. By Deligne's theorem of the
fixed part \cite{Del71}, the global sections carry a Hodge structure that
does not depend on the point, and the inclusion into each fibre is a
morphism of Hodge structures. So if $\alpha(s)$ is a Hodge class at one
point, it is a Hodge class at every point. The variational Hodge conjecture
asks for the corresponding statement about cycles: if $\alpha(s_{0})$ is
algebraic at one point $s_{0}$, then $\alpha(s)$ is algebraic at every
point $s$.

In our family, $\omega$ is a global section, it is a Hodge class
everywhere, and by Proposition~\ref{prop:densebook} it is algebraic at $s_{0}$.
The variational conjecture for $\pi\colon\mathcal A\to\mathrm{Sh}$ and this one
class says exactly that $\Sigma=\mathrm{Sh}$. Three things are known about the
conjecture in general, and none of them settles this case.
\begin{enumerate}[label=(\alph*)]
\item It holds when $p=1$, by the theorem of Lefschetz on $(1,1)$ classes.
Here $p=n\ge2$.
\item It follows from the Hodge conjecture for a smooth projective
compactification $\bar X$ of the total space. By the global invariant cycle
theorem \cite{Del71} there is a class on $\bar X$ restricting to
$\alpha(s)$, by semisimplicity of polarisable Hodge structures it can be
chosen to be a Hodge class, and the restriction of a cycle representing it
represents $\alpha(s)$ at every $s$. For our family $\bar X$ is a variety
of dimension $2n+n^{2}$, or $2n+1$ after restricting to a curve, and its
Hodge conjecture is no easier than the one we want.
\item Andr\'e's deformation principle \cite{And96} says that a flat section
which is motivated at one point is motivated at every point, and Andr\'e
also showed that every Hodge class on an abelian variety is motivated. So
$\omega(t)$ is already known to be motivated for every $t$. A motivated
class is algebraic as soon as the Lefschetz standard conjecture holds for
the varieties that enter its construction. This is the same reduction to
the standard conjectures as in Section~\ref{sec:stopsbook}, reached from the
other side.
\end{enumerate}

The rest of this section contains three results. The first changes the
question into a question about degrees. It is a variant, with Chow
varieties and arbitrary Zariski-dense sets, of the bounded criterion of
\cite{BhH8}, which is stated for Hilbert data along isogeny orbits. The
second shows that nothing built from divisors reaches $\omega$ at a member
with full Hodge group. Its proof is the argument that \cite{BhH8} gives for
Mumford's exceptional classes, carried over to the Weil family. The third
is the analogue for the tautological cycles of curves with an automorphism
of order three.

\subsection{A criterion in terms of degrees}

Measure the size of a cycle by its degree against the polarisation: for an
effective cycle $Z$ of dimension $n$ on $A_{t}$, put
$\deg_{\eta}Z=\int_{A_{t}}\cl(Z)\cup\eta^{n}$. Every cycle with rational
coefficients can be written as $\frac1m(Z^{+}-Z^{-})$ with $Z^{\pm}$
effective and $m$ a positive integer. For $N\ge1$ let $\Sigma_{N}$ be the
set of points $t$ at which
\[
  \omega(t)=\tfrac1m\bigl(\cl(Z^{+})-\cl(Z^{-})\bigr)
\]
for some effective cycles $Z^{\pm}$ on $A_{t}$ of dimension $n$ with
$\deg_{\eta}Z^{\pm}\le N$ and some integer $1\le m\le N$.

\begin{proposition}\label{prop:boundedbook}
Each $\Sigma_{N}$ is a closed algebraic subset of $\mathrm{Sh}$, and
$\Sigma=\bigcup_{N}\Sigma_{N}$. The following are equivalent.
\begin{enumerate}[label=(\roman*)]
\item $\Sigma=\mathrm{Sh}$.
\item $\Sigma_{N}=\mathrm{Sh}$ for some $N$.
\item For some $N$, the set $\Sigma_{N}$ contains a Zariski-dense subset
of $\mathrm{Sh}$.
\end{enumerate}
\end{proposition}

\begin{proof}
Effective cycles of dimension $n$ and bounded degree in the fibres of the
projective family $\mathcal A\to\mathrm{Sh}$ are parametrised by a relative Chow
variety which is of finite type and proper over $\mathrm{Sh}$ \cite[Ch.~I]{Kol96}.
Over a connected component of it, the class of the cycle is a flat section
of $R^{2n}\pi_{*}\QQ$, so the condition that
$\frac1m(\cl Z^{+}-\cl Z^{-})=\omega$ holds on whole connected components
of the fibre product of two such Chow varieties, or nowhere on them. Taking
the union over $m\le N$ of the components where it holds and projecting to
$\mathrm{Sh}$ gives $\Sigma_{N}$, which is therefore closed and algebraic. Every
point of $\Sigma$ lies in some $\Sigma_{N}$.

(i) implies (ii): $\mathrm{Sh}$ is the countable union of the closed sets
$\Sigma_{N}$, so by Baire's theorem one of them has nonempty interior, and
since $\mathrm{Sh}$ is irreducible that one is all of $\mathrm{Sh}$. (ii) implies (iii)
trivially. (iii) implies (i): a Zariski-closed set that contains a
Zariski-dense set is everything.
\end{proof}

The proposition is elementary, but it changes the shape of the question,
as \cite{BhH8} already observed in its own setting.
Proposition~\ref{prop:densebook} already gives algebraic representatives of
$\omega$ on a dense set. What is missing is a \emph{uniform bound} on their
degrees and denominators along some Zariski-dense subset. Transport along
isogenies does not give one, and it is worth seeing why.

\begin{remark}\label{rem:heckebook}
Let $\phi\colon A_{t}\to A_{t'}$ be an isogeny of degree $d$ that commutes
with $K$ and satisfies $\phi^{*}\eta'=e\,\eta$ for an integer $e\ge1$, as
the Hecke correspondences do. The two polarisations have the same type, so
comparing Euler characteristics gives
$e^{2n}\chi(\eta)=\chi(\phi^{*}\eta')=d\,\chi(\eta')=d\,\chi(\eta)$, that is
$d=e^{2n}$. For an effective cycle $Z$ of dimension $n$ the projection
formula gives
\[
  \deg_{\eta'}\phi_{*}Z=\int_{A_{t}}\cl(Z)\cup\phi^{*}\eta'^{\,n}
  =e^{n}\deg_{\eta}Z .
\]
An abelian variety has only finitely many isogenies of bounded degree, so
along any infinite subset of the Hecke orbit of $s_{0}$ the multipliers $e$
are unbounded, and so are the degrees of the transported cycles.
Proposition~\ref{prop:boundedbook} asks for \emph{other} cycles at those
points, of bounded degree. The Hecke construction gives no hint where to
find them.
\end{remark}

\subsection{A barrier at members with full Hodge group}

Fix $t\in\mathrm{Sh}^{\circ}$, write $A=A_{t}$, and let $U=U(V,H)$ be the unitary
group of the hermitian form, as an algebraic group over $\QQ$. It acts on
$V=H^{1}(A,\QQ)$, and so on $H^{*}(A,\QQ)=\bigwedge^{*}V$. Let $\mathcal C$
be the category of abelian varieties isogenous to a power of $A$, with
homomorphisms as morphisms. Choosing an isogeny $B\to A^{m}$ lets $U$ act on
$H^{1}(B,\QQ)\cong V^{m}$ diagonally. Two choices differ by an invertible
element of $\End^{0}(A^{m})$, and we shall see in a moment that these
commute with $U$, so the action does not depend on the choice.

\begin{theorem}\label{thm:barrierbook}
Let $n\ge2$ and $t\in\mathrm{Sh}^{\circ}$. Let $(\mathcal L(B))_{B\in\mathcal C}$ be
the smallest family of subspaces $\mathcal L(B)\subseteq H^{*}(B,\QQ)$
such that each $\mathcal L(B)$ is a subalgebra containing the classes of all
divisors on $B$, and $h^{*}\mathcal L(B')\subseteq\mathcal L(B)$ and
$h_{*}\mathcal L(B)\subseteq\mathcal L(B')$ for every homomorphism
$h\colon B\to B'$ in $\mathcal C$. Then every class in every $\mathcal L(B)$
is invariant under $U$. In particular no class in $\mathcal L(A)$ has the
form $\alpha+c\,\eta^{n}$ with $\alpha$ a nonzero Weil class.
\end{theorem}

\begin{proof}
\emph{Step 1: the endomorphisms.} Over $\CC$ the space $V$ splits as
$V_{+}\oplus V_{-}$, the two eigenspaces of $K$, each of
dimension $2n$. The polarisation form $E$ is invariant under the norm-one
elements of $K$, which forces $V_{+}$ and $V_{-}$ to be
isotropic and $E$ to pair them perfectly. So, as representations of
$\SU(V,H)_{\CC}\cong\SL_{2n}$, they are the standard representation and its
dual. For $2n\ge3$ these are irreducible and not isomorphic, so the
commutant of $\SU(V,H)$ in $\End_{\QQ}(V)$ has dimension two. Since the
Hodge group of $A$ is $\SU(V,H)$, this commutant is $\End^{0}(A)$, and so
$\End^{0}(A)=K$ and $\End^{0}(A^{m})=M_{m}(K)$. A homomorphism
$h\colon A^{m}\to A^{k}$ therefore acts on $H^{1}$ by a matrix with entries
in $K$, and since $U$ acts $K$-linearly on $V$, the map
$h^{*}\colon V^{k}\to V^{m}$ commutes with the diagonal action of $U$. The
same holds in $\mathcal C$ after composing with isogenies.

\emph{Step 2: the centre.} The centre of $U$ is the torus $U(1)_{K}$ of
norm-one elements of $K$, and $U=\SU(V,H)\cdot U(1)_{K}$. An element
$\lambda$ of the centre acts on $V_{+}$ by $\lambda$, through the fixed embedding of $K$, and on
$V_{-}$ by $\lambda^{-1}$, so on the summand
$\bigwedge^{a}V_{+}^{\oplus m}\otimes\bigwedge^{b}V_{-}^{\oplus m}$
of $H^{a+b}(A^{m},\CC)$ it acts by $\lambda^{a-b}$. Hence a class is
$U$-invariant exactly when it is $\SU(V,H)$-invariant and lies in the sum of
the summands with $a=b$.

\emph{Step 3: divisors.} A divisor class on $A^{m}$ is a Hodge class, and
the Hodge group of $A^{m}$ is $\SU(V,H)$ acting diagonally, so the class is
$\SU(V,H)$-invariant. Its components with $a\ne b$ lie in
$\bigwedge^{2}(V_{+}\otimes\CC^{m})$ or its conjugate. As a
representation of $\SL_{2n}$ this is a sum of copies of
$\bigwedge^{2}V_{+}$ and $\mathrm{Sym}^{2}V_{+}$, which have no
invariants when $2n\ge3$. So those components vanish, and by Step~2 the
class is $U$-invariant.

\emph{Step 4: the operations.} Cup product and $h^{*}$ commute with $U$, by
Step~1. Every element of $U$ preserves $E$, so it has determinant one on $V$
and acts trivially on the top cohomology of each $B$. So $U$ preserves the
Poincar\'e pairing, and $h_{*}$, being the adjoint of $h^{*}$ for that
pairing, commutes with $U$ as well. It follows that the family of
$U$-invariant subalgebras $H^{*}(B,\QQ)^{U}$ has all the properties in the
statement, and by minimality $\mathcal L(B)\subseteq H^{*}(B,\QQ)^{U}$.

\emph{Step 5: Weil classes.} The Weil plane, tensored with $\CC$, is
$\bigwedge^{2n}V_{+}\oplus\bigwedge^{2n}V_{-}$, the summands
with $(a,b)=(2n,0)$ and $(0,2n)$. The centre acts on them by
$\lambda^{\pm2n}$. Taking $\lambda=\tau/\bar\tau$ with $\tau$ as in
the proof of Theorem~\ref{thm:semiregbook}, this is not $1$, so the Weil plane
has no nonzero $U$-invariant vector. Since $\eta^{n}$ is $U$-invariant, a
$U$-invariant class $\alpha+c\eta^{n}$ has $\alpha$ $U$-invariant, hence
$\alpha=0$.
\end{proof}

The theorem reaches further than it looks, because Chern characters behave
well under the same operations.

\begin{corollary}\label{cor:barrierbook}
Let $t\in\mathrm{Sh}^{\circ}$, and let $F$ be an object of the derived category of
$A_{t}$ obtained from line bundles on varieties in $\mathcal C$ by finitely
many shifts, direct sums, cones, derived tensor products, and derived
pullbacks and pushforwards along homomorphisms. Then $\operatorname{ch}(F)$
is $U$-invariant, and $F$ cannot satisfy the hypothesis of
Theorem~\ref{thm:semiregbook}. This applies in particular to Fourier--Mukai
transforms whose kernel is the Mumford line bundle
$m^{*}L\otimes\mathrm{pr}_{1}^{*}L^{-1}\otimes\mathrm{pr}_{2}^{*}L^{-1}$ on
$A_{t}\times A_{t}$, and to pushforwards of line bundles along isogenies,
which by a theorem of Mukai include all simple semi-homogeneous vector
bundles \cite{Muk78}.
\end{corollary}

\begin{proof}
The Chern character is additive on exact triangles and multiplicative on
derived tensor products, it commutes with pullback, and
$\operatorname{ch}(L)=\exp c_{1}(L)$ for a line bundle. For pushforward
along a homomorphism $h$, the Grothendieck--Riemann--Roch theorem gives
$\operatorname{ch}(Rh_{*}F)=h_{*}\operatorname{ch}(F)$, because the Todd
class of an abelian variety is $1$. So $\operatorname{ch}(F)$ lies in
$\mathcal L(A_{t})$, and Theorem~\ref{thm:barrierbook} applies.
\end{proof}

So at a member with full Hodge group, nothing built from line bundles and
homomorphisms has the right Chern character. This explains a feature of
every successful construction so far. At $s_{0}$ the endomorphism algebra
is large, the corresponding group is smaller than $U$, and $\omega$ is a
polynomial in divisor classes. Markman's secant sheaves also start on
special members, built from an abelian $n$-fold and its dual, and reach the
general member only by deformation \cite{Mar25a,Mar25b}. The theorem shows
that this detour is forced: a proof has to start where the endomorphisms
are large and then deform, or it has to find at a general member a cycle
that does not come from divisors at all.

\subsection{Curves with an automorphism of order three}

The oldest source of cycles that do not come from divisors is curves. For
$K=\QQ(\sqrt{-3})$ there are two natural families of curves whose
Jacobians contain abelian varieties of Weil type, and one might hope that
cycles built from the curve reach $\omega$. The next proposition shows that
the obvious cycles do not.

Let $C$ be a smooth projective curve with an automorphism $\rho$ of order
three, and let $V\subseteq H^{1}(C,\QQ)$ be the kernel of
$1+\rho^{*}+\rho^{*2}$, the part on which $\rho$ has no fixed vectors. Then
$K$ acts on $V$ through $\rho$. Let $P\subseteq J(C)$ be the image of
$1-\rho$, so that $H^{1}(P,\QQ)$ is identified with $V$, and assume that
$P$, with the polarisation induced from $J(C)$, is of Weil type for $K$.
Let $U=U(V,H)$ act on $H^{1}(C,\QQ)=V\oplus H^{1}(C,\QQ)^{\rho}$, trivially
on the second summand, and trivially on $H^{0}$ and $H^{2}$.

\begin{proposition}\label{prop:curvesbook}
Let $R^{*}(C^{m})\subseteq H^{*}(C^{m},\QQ)$ be the subring generated by
the classes of the diagonals $\{x_{i}=x_{j}\}$, the graphs
$\{x_{j}=\rho^{k}(x_{i})\}$ for $k=1,2$, and the pullbacks of the point
class from each factor. Let $f\colon C^{m}\to P$ be a morphism of the form
$h\circ(\mathrm{AJ}\times\dots\times\mathrm{AJ})$, where $\mathrm{AJ}$ is an
Abel--Jacobi map $C\to J(C)$ and
$h(y_{1},\dots,y_{m})=(1-\rho)\sum_{i}\psi_{i}(y_{i})$ with
$\psi_{i}\in\ZZ[\rho]$. Then every class in $f_{*}R^{*}(C^{m})$ is
$U$-invariant, and in particular none has the form $\alpha+c\,\eta^{n}$
with $\alpha$ a nonzero Weil class.
\end{proposition}

\begin{proof}
The summand $V$ is orthogonal to $H^{1}(C,\QQ)^{\rho}$ under cup product,
because the pairing is $\rho$-invariant and $\rho$ acts on the two summands
with no common eigenvalue. The group $U$ preserves the pairing on $V$, which
is the polarisation form of $P$ up to a constant. So the action of $U$ on
$H^{1}(C,\QQ)$ is symplectic, it extends to an action on $H^{*}(C^{m},\QQ)$
by algebra automorphisms, and it fixes the top class.

The class of the graph of an endomorphism $T$ of $H^{1}$, in the
$H^{1}\otimes H^{1}$ part of $H^{2}(C\times C)$, is the tensor
corresponding to $T$ under Poincar\'e duality. It is invariant under every
symplectic $g$ that commutes with $T$. The diagonal is the case $T=1$, and
the graphs of $\rho^{k}$ are the cases $T=\rho^{\pm k}$, which commute with
$U$ because $U$ is $K$-linear. The other components of these classes, and
the point classes, lie in $H^{0}\otimes H^{2}$ and $H^{2}\otimes H^{0}$,
where $U$ acts trivially. So $R^{*}(C^{m})$ is $U$-invariant.

The pullback $f^{*}\colon H^{1}(P)=V\to H^{1}(C)^{\oplus m}$ is given in
each coordinate by an element of $\ZZ[\rho]$, that is, by multiplication by
an element of $K$. So it lands in $V^{\oplus m}$ and commutes with $U$,
which acts $K$-linearly. (This is where the special form of $h$ is used: a
general $K$-linear map of $V$ does not commute with $U$.) Hence $f^{*}$ commutes
with $U$ on all of $H^{*}(P)$, and $f_{*}$, its adjoint for the Poincar\'e
pairings, does too. The last assertion follows as in Step~5 of the proof
of Theorem~\ref{thm:barrierbook}.
\end{proof}

Two families fit the hypotheses. In the first, $C$ is the cyclic triple
cover $y^{3}=\prod_{i=1}^{n+1}(x-a_{i})\prod_{j=1}^{n+1}(x-b_{j})^{2}$ of
the projective line. By Riemann--Hurwitz it has genus $2n$, and by the
Chevalley--Weil formula the two nontrivial eigenspaces of $\rho$ on
holomorphic differentials both have dimension $n$, so $P=J(C)$ is of Weil
type. These curves depend on $2n-1$ parameters, which is less than $n^{2}$
for $n\ge2$, so their Jacobians fill only a proper subvariety of $\mathrm{Sh}$. In
the second, $C\to C'$ is an \'etale cyclic triple cover of a curve of genus
$n+1$. Then $C$ has genus $3n+1$, the two nontrivial eigenspaces both have
dimension $n$, and $P$ is the Prym variety, of dimension $2n$. These
depend on $3n$ parameters, and for $n\le3$ that is at least $n^{2}$. Here
much more is known than the proposition might suggest. Schoen proved that
the Weil classes of these Prym varieties are algebraic \cite{Sch88}, so for
$n\le3$ the Pryms give whole families of Weil type in which the conjecture
holds, and for $n\ge4$ they give a locus of dimension $3n$ inside $\mathrm{Sh}$ on
which $\omega$ is algebraic and whose very general point has full Hodge
group. The proposition says something about how this can happen. Schoen's
cycles represent a nonzero Weil class, so they cannot be the image of a
tautological class under a map of the form in the proposition. Whatever
they are built from, it is a cycle on a power of $C$, or a map, outside the
tautological world. That is exactly the kind of source that
Theorem~\ref{thm:barrierbook} says a proof needs, and it is the most concrete
lead this chapter can offer (see Section~\ref{sec:leftbook}).

\section{What is left}\label{sec:leftbook}

Put together, the three results narrow the search. At a member with full
Hodge group, $\omega$ cannot be reached by anything built from divisors and
homomorphisms (Theorem~\ref{thm:barrierbook} and
Corollary~\ref{cor:barrierbook}), nor by the tautological cycles of a curve
with an automorphism of order three (Proposition~\ref{prop:curvesbook}). At
special members it can be reached, and Hecke transport spreads those
representatives over a dense set, but with degrees that the transport
itself drives to infinity, and Proposition~\ref{prop:boundedbook} says that
bounded degree on a Zariski-dense set is exactly what is needed. Any proof
therefore has to supply one of three things: a uniform bound on degrees
along a Zariski-dense set of special members; a deformation argument that
carries a cycle from a special member to the general one, which is the
semiregularity route of Theorem~\ref{thm:semiregbook}; or a source of cycles at
a general member that is neither divisorial nor tautological. The third
exists for $\QQ(\sqrt{-3})$: Schoen's cycles on Prym varieties. Their
locus has dimension $3n$, so for $n\ge4$ they reach only a thin part of
$\mathrm{Sh}$. The obvious question is whether one of Schoen's cycles is
semiregular. The construction makes the question concrete. On the
symmetric power $C'^{(2n)}$ of the base curve, the canonical system
$|K_{C'}|\cong\PP^{n}$ is the fibre of the Abel--Jacobi map over the
canonical class. The triple cover defines an \'etale cover of $C'^{(2n)}$
of degree three, and over $|K_{C'}|$, which is simply connected, it splits
into three copies of $\PP^{n}$. The differences of the copies, weighted by
the characters of order three, are Schoen's cycles
\cite[Theorem~2.0 with $r=0$]{Sch88}, and Patel and Zhang reach the same
copies by base change along an isogeny onto the Jacobian
\cite[Theorem~1.1]{PZ25}. On the Prym the class of one copy is a Weil
class plus a multiple of $\eta^{n}$. When
that multiple is nonzero, this is exactly the shape that Theorem~\ref{thm:semiregbook}
allows and that the obstruction from Weil tori does not exclude. If one of them were
semiregular, Theorem~\ref{thm:semiregbook} would give $\Sigma=\mathrm{Sh}$ for the
family of that dimension. The construction also suggests why this may
fail. The Abel--Jacobi curve $C\subset J(C)$ comes from the same map, and
its class $\theta^{g-1}/(g-1)!$ is algebraic on every principally polarized
abelian variety; yet by the criterion of Matsusaka and Ran a curve with this
class exists only on Jacobians and products of Jacobians
\cite{Mat59,Ran81}. For $g\ge4$ these fill a proper subvariety of the
moduli space, so by Bloch's theorem the Abel--Jacobi curve is not
semiregular. Schoen's cycles are the analogue one step up, and nothing in
their construction moves them off the Prym locus. This is an analogy, not a
proof, and I have not been able to decide the question. Outside
$\QQ(\sqrt{-3})$, none of the three is available for $n\ge3$ in general.

To sum up what this chapter proves. The algebraic locus of the Weil class
is a countable union of closed algebraic subsets, it is dense, and it is
either everything or meagre. A positive-dimensional stratum of it through
a member with full Hodge group has full monodromy. A single semiregular
cycle or complex, at any member, with class a nonzero Weil class plus a
multiple of $\eta^{n}$, would close the locus, while an object whose Chern
character is a pure Weil class is never semiregular. The problem is the
variational Hodge conjecture for one class. It is equivalent to a bound on
degrees along a Zariski-dense set, and at a member with full Hodge group it
cannot be solved with divisors or with the tautological cycles of curves.

What it does not prove is the conclusion. For $n\ge4$, and for $n=3$
outside the known families, the Hodge conjecture for abelian varieties of
Weil type remains open. Even a complete solution would leave the other
exceptional classes on abelian varieties, and the conjecture for varieties
that are not abelian, which is the conjecture in full.
